% !TeX program = pdflatex
% ─────────────────────────────────────────────────────────────────────────────
%  papers/combinacao.tex — A FRONTEIRA DA COMBINAÇÃO
%
%  Regista a família experimental C0–C8 (lib/arena_combinacao.mjs) sem definir C.
%  Continuação de papers/redes.tex (§RG6 FECHADO). C0--L7 intacto; M1--M2 e E_∂
%  registados. Não promove lexMax. P6 TRAVADA. Compila:
%    cd papers && pdflatex combinacao.tex
% ─────────────────────────────────────────────────────────────────────────────
\documentclass[11pt,a4paper]{article}
\usepackage[utf8]{inputenc}\usepackage[T1]{fontenc}\usepackage[brazil]{babel}
\usepackage{amsmath,amssymb,amsthm}
\usepackage[margin=2.4cm]{geometry}
\usepackage{booktabs}\usepackage{array}\usepackage{xcolor}
\usepackage[htt]{hyphenat}
\usepackage[hidelinks,breaklinks]{hyperref}

\sloppy
\emergencystretch=2em

\definecolor{cinza}{gray}{0.35}
\newcommand{\code}[1]{\texttt{\small #1}}
\newcommand{\caixa}[1]{%
  \par\smallskip\noindent
  \begin{center}
  \fbox{\begin{minipage}{0.94\textwidth}\centering\smallskip #1\smallskip\end{minipage}}
  \end{center}\smallskip}
\newcommand{\abs}[1]{\lvert #1\rvert}
\newcommand{\medido}[1]{\par\medskip\noindent{\small\color{cinza}\textbf{Medido.} #1}\par\medskip}
\newcommand{\fis}[2]{\textbf{#2}~\textnormal{(\code{#1})}}
\newcommand{\redes}[1]{\code{redes.tex} §\textit{#1}}
\newcommand{\sepcol}{\quad\penalty0$\big|$\quad\penalty0}
\providecommand{\interfacen}[1]{}
\newtheorem{teorema}{Teorema}
\newtheorem{lema}[teorema]{Lema}
\newtheorem{definicao}[teorema]{Definição}
\newtheorem{corolario}[teorema]{Corolário}
\newtheorem{proposicao}[teorema]{Proposição}
\newtheorem{hipotese}[teorema]{Hipótese experimental}
\renewcommand{\proofname}{Demonstração}

% ─────────────────────────────────────────────────────────────────────────────
%  papers/gkcapa.tex — a capa do Reino, mínima, para os papers em `article`.
%
%  O estilo.tex completo é do `report` (títulos de capítulo, skak, …). Aqui fica
%  só o que a capa precisa, para o paper continuar a compilar sozinho com
%  pdflatex. O tradutor próprio (tests/tex) carrega o estilo.tex da raiz / o
%  symlink papers/estilo.tex e avalia o mesmo `\gkcapa`.
% ─────────────────────────────────────────────────────────────────────────────
\definecolor{ouro}{HTML}{B8912F}
\definecolor{ouroclaro}{HTML}{F3E9CE}
\definecolor{tinta}{HTML}{1E1B16}
\definecolor{regua}{HTML}{6E6858}
\providecommand{\gknota}{\fontsize{7.62}{11.05}\selectfont}
\providecommand{\gktexto}{\fontsize{10.50}{15.22}\selectfont}
\providecommand{\gksub}{\fontsize{12.33}{17.87}\selectfont}
\providecommand{\gksec}{\fontsize{14.47}{20.98}\selectfont}
\providecommand{\gkcap}{\fontsize{16.99}{24.63}\selectfont}
\providecommand{\gktit}{\fontsize{23.42}{33.95}\selectfont}
\providecommand{\Prj}{\mathbb{P}}
\providecommand{\Trf}{\mathcal{T}}
\providecommand{\gkcapa}[3]{%
\title{\vspace{-0.6cm}
{\color{ouro}\rule{\textwidth}{1.4pt}}\\[6mm]
\textbf{\gktit\color{tinta} Reino Dourado}\\[3mm]{\gksec\itshape\color{regua} Golden Kingdom}\\[8mm]
{\gkcap\scshape\color{ouro} #1}\\[5mm]
{\gksub\color{regua} #2}\\[2mm]
{\gktexto\color{regua}\itshape #3}\\[7mm]
{\color{ouro}\rule{\textwidth}{0.6pt}}\\[9mm]{\gknota\color{regua}\begin{minipage}{\textwidth}\setlength{\parindent}{0pt}%
\textbf{Aviso de Isenção Algorítmica:} O universo, as leis e as entidades contidas nestas páginas ---
incluindo felinos matriciais, aves de rapina algébricas e amuletos de controle de fluxo --- são
construções de pura ficção formal. Esta obra não descreve a realidade física, não serve como
engenharia reversa do cosmos e não constitui doutrina religiosa. Qualquer semelhança com bugs reais do seu
sistema operacional é mera coincidência (ou falta de reza). Prossiga por sua própria conta e
risco.\end{minipage}}}
\author{}\date{}}


\begin{document}
\interfacen{7}
\gkcapa{A Fronteira da Combinação}
       {overwrite, cópia, acumulação e a classe admissível de combinações}
       {C0--L7: a máquina encontrou uma \emph{classe}, não uma lei única}
\maketitle

\begin{abstract}\noindent
\textbf{De onde parte.} De \code{papers/redes.tex} §RG6: a arquitectura host determina
\textbf{overwrite} quando múltiplos focos incidem na mesma representação, mas \textbf{não}
determina uma lei $C(x_1,x_2)$ entre incidências concorrentes. Este documento regista a
escada experimental \textbf{C0--L7} (\code{lib/arena\_combinacao.mjs}) --- do que o suporte
\emph{faz} até à classe de combinações que a infraestrutura \emph{admite}, sem postular $C$.

\medskip\noindent\textbf{Tese final.}
\[
\boxed{\text{a máquina encontrou uma \textbf{classe} de combinações,
não uma lei única de combinação.}}
\]
Existência de $f\colon X\times X\to X$ bem definido e bilateralmente dependente:
\textbf{sim} (D0--D3). Unicidade: \textbf{não} (K0--K7). A incidência $\zeta/\mu/S$ determina
acumulação/deconvolução, mas \textbf{não} determina por si só a lei de combinação (L7).

\medskip\noindent\textbf{Verificação.}
\code{tests/redes\_combinacao.js} (231 asserts);
\code{tools/auditoria\_rg6.bat} (bateria completa, após §RG6).
M1--M2 e $E_{\partial}$ registados; $C$ continua reservada.
\end{abstract}

% ═══════════════════════════════════════════════════════════════════════════
\section{A regra deste documento}\label{comb:regra}

\noindent\textbf{Princípio.}
\caixa{%
  não procurar a lei que queremos;\\
  procurar a lei que o suporte permite.%
}

\medskip\noindent\textbf{Proibições explícitas.}
\begin{enumerate}\itemsep3pt
\item \textbf{Não definir} $C\colon X\times X\to X$ neste paper;
\item \textbf{Não rotular} overwrite como interferência, Hebb, soma ou Hopfield;
\item \textbf{Não confundir} $G_{\mathrm{event}}$ com $G_{\mathrm{state}}^{\mathrm{surv}}$
(\redes{RG6});
\item \textbf{Não confundir} acumulação $A$ (dois slots físicos) com combinação $C$
(um slot, dois payloads);
\item \textbf{Não alterar} \code{redes.tex} --- cumpriu o papel de interface;
\item \textbf{Não identificar} L5/K6 com $\tau$ da cisão nem com $D_{\mathrm{can}}$;
\item \textbf{Não abrir P6} --- reservada para comparação futura (§\ref{comb:hebb}).
\end{enumerate}

\medskip\noindent\textbf{O que se mede.} Protocolos nomeados C0--C8 sobre a arena
multifocal (\code{lib/arena\_multifocal.mjs}): quatro focos, célula lógica $X$,
slots \code{OFF\_IN}/\code{OFF\_OUT}, registo de eventos, overwrites e métricas $G$.
Cada protocolo responde uma pergunta \emph{sem} escolher a lei de antemão.

\medskip\noindent\textbf{Escada completa (congelada).}
{\small
\[
\begin{array}{c}
\text{C0--C8}\ \to\ \text{caracterização}\ \to\ \text{operações observáveis}\\
\text{S0--S4}\ \to\ \text{capacidade do suporte}\ \to\ \text{preservação}\\
\text{R0--R4}\ \to\ \text{redução deliberada}\ \to\ R\colon(x_1,x_2)\mapsto x\\
\text{D0--D3}\ \to\ \text{sonda do cruzado}\ \to\ \text{dependência bilateral}\\
\text{K0--K7}\ \to\ \text{caracterização}\ \to\ \text{não-unicidade}\\
\text{L0--L7}\ \to\ \text{invariantes}\ \to\ \mathcal A_{\min}\ \text{e discriminantes}\\
\text{M1--M2}\ \to\ \text{ordem lex}\ \to\ \text{restringe a classe, não escolhe }C\\
\text{$E_{\partial}$}\ \to\ \text{covariância sob }D_{\mathrm{can}}\ \to\ \text{não selecciona }C
\end{array}
\]}
Este documento fecha na \textbf{classificação}; $C(x,y)$ permanece reservada.

% ═══════════════════════════════════════════════════════════════════════════
\section{Objeto e cinco leituras de $G$}\label{comb:objeto}

\noindent\textbf{Objecto experimental.} Arena única de $65536$ bytes; quatro diários
$I_a$ (focos $F_1,\ldots,F_4$) incidem sobre a célula lógica $X$ via slots físicos
\code{OFF\_IN} e \code{OFF\_OUT}. Cada incidência regista foco, operação, payload e
slot --- \textbf{sem} chamar \code{interpretar.c}.

\medskip\noindent\textbf{Cinco leituras de $G$} (herdadas de \redes{RG6}):
{\small
\begin{align*}
G_{\mathrm{event}}(x) &= \text{número de eventos/visitas em }x;\\
G_{\mathrm{focus}}(x) &= \text{focos distintos que incidiram sobre }x;\\
G_{\mathrm{state}}^{\mathrm{esc}}(x) &= \text{payloads distintos \emph{escritos} em }x;\\
G_{\mathrm{state}}^{\mathrm{surv}}(x) &= \text{payloads distintos \emph{sobreviventes}
no suporte físico de }x;\\
G_{\mathrm{slots}}^{\mathrm{surv}}(x) &= \text{slots físicos não vazios ligados a }x.
\end{align*}}
Implementação: \code{gEvent}, \code{gFocus}, \code{gStateWritten},
\code{gStateSurviving}, \code{gSlotsSurviving} em \code{lib/arena\_combinacao.mjs}.

\medskip\noindent\textbf{Distinção central (já congelada em C0--C4).}
\caixa{%
  $\text{overwrite}\neq\text{cópia}\neq\text{acumulação}\neq\text{combinação}$%
}
Fronteira experimental limpa:
\[
\boxed{C2\colon\ G_{\mathrm{state}}^{\mathrm{surv}}=2
\qquad\text{vs}\qquad
C3\colon\ G_{\mathrm{state}}^{\mathrm{surv}}=1.}
\]

% ═══════════════════════════════════════════════════════════════════════════
\section{C0 --- overwrite}\label{comb:c0}

\noindent\textbf{Protocolo.} Estado inicial $X_0$; $F_1\!\to\!X\colon x_1$;
$F_2\!\to\!X\colon x_2$ (mesmo slot \code{OFF\_IN}); $F_3$ lê.

\medskip\noindent\textbf{Resultado observado.}
\[
C_{\mathrm{ow}}(x_1,x_2)=x_2.
\]
Registo \code{overwrite\_single\_slot}; $G_{\mathrm{state}}^{\mathrm{surv}}(X)=1$;
$G_{\mathrm{state}}^{\mathrm{esc}}(X)=2$.

\medskip\noindent\textbf{Leitura.} Overwrite é \textbf{política do suporte} --- uma
coordenada por slot; segunda escrita substitui (\fis{fis:def:incid}{incidência}).
Não implica lei de combinação.

% ═══════════════════════════════════════════════════════════════════════════
\section{C1 --- duplicação}\label{comb:c1}

\noindent\textbf{Protocolo.} $F_1$ escreve $x_1$ em \code{OFF\_IN}; $F_2$ escreve
$x_1$ em \code{OFF\_OUT}; $F_3,F_4$ leem.

\medskip\noindent\textbf{Resultado observado.} Dois slots físicos ocupados;
\textbf{um} payload distinto sobrevive ($G_{\mathrm{state}}^{\mathrm{surv}}=1$,
$G_{\mathrm{slots}}^{\mathrm{surv}}=2$). Sem overwrite.

\medskip\noindent\textbf{Leitura.} \textbf{Cópia}, não duas realizações independentes.
A mesma informação numérica aparece em duas coordenadas físicas da mesma célula lógica.

% ═══════════════════════════════════════════════════════════════════════════
\section{C2 --- acumulação}\label{comb:c2}

\noindent\textbf{Protocolo.} $F_1\!\to\!\code{OFF\_IN}\colon x_1$;
$F_2\!\to\!\code{OFF\_OUT}\colon x_2$; leituras.

\medskip\noindent\textbf{Resultado observado.}
\[
A(x_1,x_2)=(x_1,x_2).
\]
$G_{\mathrm{state}}^{\mathrm{surv}}(X)=2$; $G_{\mathrm{slots}}^{\mathrm{surv}}=2$;
sem overwrite.

\medskip\noindent\textbf{Leitura.} Acumulação usa \textbf{dois suportes físicos}
($V=\mathbb{Z}_c^X$ --- duas coordenadas, mesma célula lógica). Distinto de
$C(x_1,x_2)=x_3$ (§\ref{comb:c3}).

% ═══════════════════════════════════════════════════════════════════════════
\section{C3 --- tentativa de combinação}\label{comb:c3}

\noindent\textbf{Protocolo.} $F_1\!\to\!X\colon x_1$; $F_2\!\to\!X\colon x_2$
(mesmo slot); $F_3$ lê.

\medskip\noindent\textbf{Resultado observado.} Estado final $x_2$ (overwrite);
\textbf{não} existe $x_3$; \code{leiC: null}.

\medskip\noindent\textbf{Leitura.} Slot único com payloads distintos \textbf{não produz}
combinação algébrica --- reproduz overwrite. A pergunta $C(x_1,x_2)=?$ permanece aberta.

% ═══════════════════════════════════════════════════════════════════════════
\section{C4 --- recuperação}\label{comb:c4}

\noindent\textbf{Protocolo.} Executa C2; relê o par acumulado.

\medskip\noindent\textbf{Resultado observado.} Par $\{in\colon x_1,\ out\colon x_2\}$
recuperado; \code{combinado: null}. Acumulação preservada ($G_{\mathrm{state}}^{\mathrm{surv}}=2$).

\medskip\noindent\textbf{Leitura.} Recuperação lê o par --- não funde incidências
(\fis{fis:thm:zetamu}{ciclo $\zeta$; recuperação $\mu$}). $\mu$ sobre série finita
reservada para experimentos posteriores.

% ═══════════════════════════════════════════════════════════════════════════
\section{C5 --- ordem}\label{comb:c5}

\noindent\textbf{Protocolo.} Mesmos payloads, slots e focos; apenas a ordem temporal
muda:
\begin{align*}
A_{12}(x_1,x_2)&\colon\ F_1\!\to\!\code{IN}(x_1),\ F_2\!\to\!\code{OUT}(x_2);\\
A_{21}(x_2,x_1)&\colon\ F_2\!\to\!\code{OUT}(x_2),\ F_1\!\to\!\code{IN}(x_1).
\end{align*}

\medskip\noindent\textbf{Resultado observado.}
\begin{center}\footnotesize
\begin{tabular}{@{}lcc@{}}
\toprule
medida & $A_{12}$ & $A_{21}$ \\
\midrule
estado final & $(x_1,x_2)$ & $(x_1,x_2)$ \\
$G_{\mathrm{event}}$ & $2$ & $2$ \\
$G_{\mathrm{focus}}$ & $2$ & $2$ \\
$G_{\mathrm{state}}^{\mathrm{esc}}$ & $2$ & $2$ \\
$G_{\mathrm{state}}^{\mathrm{surv}}$ & $2$ & $2$ \\
timeline & \code{F1→IN; F2→OUT} & \code{F2→OUT; F1→IN} \\
\bottomrule
\end{tabular}
\end{center}

\medskip\noindent\textbf{Resultado congelado.}
\[
\boxed{A_{12}=A_{21}\ \text{como estado};\quad\text{timeline}\neq.}
\]

\medskip\noindent\textbf{Leitura.} A acumulação de \textbf{duas} incidências em slots
distintos é \textbf{comutativa como estado}. A ordem temporal distingue-se na
realização (eventos), não no estado final.

% ═══════════════════════════════════════════════════════════════════════════
\section{C6 --- repetição}\label{comb:c6}

\noindent\textbf{Protocolo.} Compara $(x_1,x_1)$ com $(x_1,x_2)$ em duas variantes:
\begin{enumerate}\itemsep2pt
\item \textbf{dual} --- $F_1\!\to\!\code{IN}(x_1)$, $F_2\!\to\!\code{OUT}(x_1)$;
\item \textbf{slot único} --- $F_1\!\to\!X\colon x_1$, $F_2\!\to\!X\colon x_1$.
\end{enumerate}

\medskip\noindent\textbf{Resultado observado.}
\begin{center}\footnotesize
\begin{tabular}{@{}lccc@{}}
\toprule
variante & estado & $G_{\mathrm{state}}^{\mathrm{surv}}$ & overwrite? \\
\midrule
dual $(x_1,x_1)$ & $(x_1,x_1)$ & $1$ & não \\
slot $(x_1,x_1)$ & $x_1$ & $1$ & não (payload idêntico) \\
C2 $(x_1,x_2)$ & $(x_1,x_2)$ & $2$ & --- \\
\bottomrule
\end{tabular}
\end{center}

\medskip\noindent\textbf{Leitura.} $A(x,x)$ dual produz \textbf{cópia} (dois slots,
um payload). No slot único, comportamento \textbf{idempotente} sem registo de overwrite
(política: overwrite só quando payloads \emph{distintos}).

% ═══════════════════════════════════════════════════════════════════════════
\section{C7 --- três incidências}\label{comb:c7}

\noindent\textbf{Protocolo.} Terceira incidência sobre o par acumulado; compara
dois aninhamentos no suporte dual:
\begin{align*}
\text{esquerda:}\ &A\bigl(A(x_1,x_2),x_3\bigr)
\ \colon\ F_1\!\to\!\code{IN}(x_1),\ F_2\!\to\!\code{OUT}(x_2),\ F_3\!\to\!\code{IN}(x_3);\\
\text{direita:}\ &A\bigl(x_1,A(x_2,x_3)\bigr)
\ \colon\ F_2\!\to\!\code{OUT}(x_2),\ F_3\!\to\!\code{IN}(x_3),\ F_1\!\to\!\code{IN}(x_1).
\end{align*}

\medskip\noindent\textbf{Resultado observado.}
\[
A\bigl(A(x_1,x_2),x_3\bigr)=(x_3,x_2),
\qquad
A\bigl(x_1,A(x_2,x_3)\bigr)=(x_1,x_2).
\]
Ambos registam overwrite em \code{OFF\_IN} na terceira incidência; estados finais
\textbf{distintos}.

\medskip\noindent\textbf{Resultado congelado --- correção conceptual.}
\[
\boxed{
A_{12}=A_{21}
\quad\text{mas}\quad
A\bigl(A(x_1,x_2),x_3\bigr)\neq A\bigl(x_1,A(x_2,x_3)\bigr).
}
\]
\textbf{Comutatividade em duas incidências não implica associatividade em três.}

\medskip\noindent\textbf{Leitura.} C5 mostrou comutatividade \emph{como estado} para
duas incidências em slots distintos. C7 mostra que o \textbf{suporte finito com
overwrite} quebra a composição de três incidências: a terceira incidência em
\code{OFF\_IN} sobrescreve uma coordenada do par, e \emph{qual} coordenatura
sobrevive depende da ordem e do aninhamento interpretado.

\medskip\noindent\textbf{Consequência.} A representação física \textbf{não} é detalhe
de implementação. A composição depende de \emph{onde} a incidência chega e de
\emph{qual} célula física ela sobrescreve.

% ═══════════════════════════════════════════════════════════════════════════
\section{C8 --- colisão}\label{comb:c8}

\noindent\textbf{Protocolo.} Situação mínima em que $C(x_1,x_2)=?$ torna-se
inevitável:
\[
F_1\!\to\!X\colon x_1,\qquad F_2\!\to\!X\colon x_2,
\]
duas fontes, mesma célula lógica, \textbf{mesmo slot físico}, payloads distintos.

\medskip\noindent\textbf{Resultado observado.} Estado final $x_2$; overwrite
registado; $G_{\mathrm{state}}^{\mathrm{surv}}=1$; \code{leiC: null}.

\medskip\noindent\textbf{Leitura.} O objectivo de C8 \textbf{não} é descobrir o
resultado (já conhecido como overwrite desde C0/C3), mas medir as \textbf{condições
de colisão}:
\caixa{%
  2 focos $\cdot$ mesma célula $X$ $\cdot$ mesmo slot $\cdot$ payloads distintos%
}

% ═══════════════════════════════════════════════════════════════════════════
\section{O que já foi demonstrado}\label{comb:demonstrado}

\noindent\textbf{Tabela C0--C8.}
\begin{center}\footnotesize
\begin{tabular}{@{}>{\raggedright\arraybackslash}p{0.38\textwidth}>{\raggedright\arraybackslash}p{0.52\textwidth}@{}}
\toprule
propriedade & resultado \\
\midrule
overwrite & observado (C0, C3, C8) \\
cópia & observado (C1, C6 dual) \\
acumulação & observado (C2) \\
combinação & \textbf{não determinada} (C3, C8: \code{leiC: null}) \\
recuperação do par & observado (C4) \\
comutatividade de 2 incidências & \textbf{sim, como estado} (C5) \\
associatividade de 3 incidências & \textbf{não} (C7; limitação do suporte cap~2) \\
$A(x,x)$ & cópia (dual) / idempotente no slot único (C6) \\
colisão mínima & 2 focos + mesmo $X$ + mesmo slot + payloads distintos (C8) \\
\bottomrule
\end{tabular}
\end{center}

\medskip\noindent\textbf{S0--S4 (capacidade).}
$\text{perda por overwrite}\neq\text{lei de composição}$. Capacidade~2: $G_{\mathrm{state}}^{\mathrm{surv}}=2$
(S1); terceira incidência perde um payload (S2). Capacidade~3: três payloads recuperáveis (S3);
não-associatividade de C7 era \textbf{limitação física}, não propriedade intrínseca.

\medskip\noindent\textbf{R0--R4 (redução).} Redução deliberada $R\colon(x_1,x_2)\mapsto x$ com
capacidade $\ge2$: preservar (R0), seleccionar $x_2$ (R1), seleccionar $x_1$ (R2),
codificar $x_1|x_2$ (R3), descartar (R4). Controle causal:
$R_{\mathrm{surv}}(\mathrm{S0})=R_{\mathrm{surv}}(\mathrm{R1})$ mas
$\mathrm{real}(\mathrm{S0})\neq\mathrm{real}(\mathrm{R1})$ --- overwrite acidental
$\neq$ redução deliberada.

% ═══════════════════════════════════════════════════════════════════════════
\section{D0--D3: sonda do cruzado}\label{comb:d}

\noindent\textbf{Metodologia.} Não se postula distributividade antes da medição.
Família de sondas sobre suporte preservador:
\begin{itemize}\itemsep2pt
\item D0 --- controle: $R_\theta=\mathrm{id}$, baseline;
\item D1 --- fixa $x_2$, varia $x_1\to x_1'$, mede $\Delta x$;
\item D2 --- fixa $x_1$, varia $x_2\to x_2'$, mede $\Delta x$;
\item D3 --- matriz $2\times2$; classifica seleção, codificação ou candidato.
\end{itemize}

\medskip\noindent\textbf{Resultado D3 (bateria de reduções).}
\begin{center}\footnotesize
\begin{tabular}{@{}ll@{}}
\toprule
$R$ & leitura \\
\midrule
keep1 & seleção de $x_1$ \\
keep2 & seleção de $x_2$ \\
encode & codificação (R3-like) \\
discard & constante $\varnothing$ \\
lexMax & \textbf{dependência dupla} --- candidato \\
rectCell & \textbf{dependência dupla} --- candidato \\
\bottomrule
\end{tabular}
\end{center}

\medskip\noindent\textbf{Primeira existência medida.}
\[
\boxed{\exists f\colon X\times X\to X\ \text{com dependência observável nos dois argumentos.}}
\]
Ainda \textbf{não} $f=C$; ainda menos ``distributividade''.

% ═══════════════════════════════════════════════════════════════════════════
\section{K0--K7: não-unicidade}\label{comb:k}

\noindent\textbf{Caracterização} dos candidatos lexMax e rectCell:
\begin{center}\footnotesize
\begin{tabular}{@{}lcc@{}}
\toprule
propriedade & lexMax & rectCell \\
\midrule
$f(x,y)=f(y,x)$ & sim & não \\
$f(f(x,y),z)=f(x,f(y,z))$ & sim & não \\
$f(y,y)=y$ & sim & não \\
bem definida (mesmo par $\Rightarrow$ mesmo $f$) & sim & sim \\
\bottomrule
\end{tabular}
\end{center}

\medskip\noindent\textbf{K7.} lexMax $\neq$ rectCell em todos os pares testados ---
\textbf{duas leis}, não duas representações da mesma.

\medskip\noindent\textbf{Resultado congelado.}
\[
\boxed{\text{dependência bilateral é \textbf{necessária}, não \textbf{suficiente}
para determinar $f$.}}
\]
\[
\boxed{\text{existência de combinação}\ \neq\ \text{unicidade da combinação.}}
\]

% ═══════════════════════════════════════════════════════════════════════════
\section{L0--L7: classe admissível e invariantes discriminantes}\label{comb:l}

\noindent\textbf{Classe mínima observada.}
\[
\mathcal A_{\min}
=\bigl\{f\colon X\times X\to X:
\text{bilateral, bem definido, }G\text{- e }\zeta/\mu\text{-compatível,
K6-invariante}\bigr\}.
\]
Implementação: L0--L7 em \code{lib/arena\_combinacao.mjs}; $\zeta,\mu$ conforme
\code{catalogo.tex} / \code{lib/incidencia.h} ($S$ gera $\zeta=\sum S^j$, $\mu=1-S$,
$\mu\zeta=\mathrm{id}$).
\emph{L5 não é $\tau$}: é a invariância de K6 (mesmo par, ordem de escrita
permutada $\Rightarrow$ mesmo resultado). Não é a cisão do par
(\code{fis:obs:cisao}) nem o $D_{\mathrm{can}}$ de \code{fis:thm:troca-realizacao}.

\medskip\noindent\textbf{Tabela L0--L7.}
\begin{center}\footnotesize
\begin{tabular}{@{}lcc@{}}
\toprule
invariante & lexMax & rectCell \\
\midrule
L0 dependência bilateral & $\checkmark$ & $\checkmark$ \\
L1 bem-definição & $\checkmark$ & $\checkmark$ \\
L2 comutatividade & $\checkmark$ & $\times$ \\
L3 idempotência & $\checkmark$ & $\times$ \\
L4 associatividade & $\checkmark$ & $\times$ \\
L5 invariância K6 (ordem de escrita) & $\checkmark$ & $\checkmark$ \\
L6 compatibilidade $G$ & $\checkmark$ & $\checkmark$ \\
L7 compatibilidade $\zeta/\mu$ & $\checkmark$ & $\checkmark$ \\
\bottomrule
\end{tabular}
\end{center}

\medskip\noindent\textbf{Membros de $\mathcal A_{\min}$.}
$\operatorname{rectCell}\in\mathcal A_{\min}$ e $\operatorname{lexMax}\in\mathcal A_{\min}$.
Acrescentando $\{\text{comutatividade, idempotência, associatividade}\}$, a classe encolhe
e rectCell sai; lexMax sobrevive --- \textbf{classificação}, não escolha de $C$.

\medskip\noindent\textbf{Conclusão sobre $\zeta/\mu$.}
\[
\boxed{\text{a incidência determina acumulação/deconvolução,
mas não determina por si só a lei de combinação.}}
\]
A ponte $\zeta+\mu\Longrightarrow C$ \textbf{não} é suficiente; medido empiricamente (L7).

\medskip\noindent\textbf{Papel dos candidatos.}
\begin{itemize}\itemsep2pt
\item \textbf{lexMax} --- sobrevivente sob exigência de semilattice; \textbf{não} declarado $C$;
\item \textbf{rectCell} --- contraexemplo à suficiência da dependência bilateral.
\end{itemize}

\medskip\noindent\textbf{Princípio metodológico.}
\caixa{%
  a máquina não recebeu a lei;\\
  recebeu critérios para procurar uma realização.%
}

% ═══════════════════════════════════════════════════════════════════════════
\section{M1--M2 e $E_{\partial}$: duas autoridades, nenhuma lei}\label{comb:me}

\noindent\textbf{Núcleo C0--L7 intocado.} Depois de L7 mediram-se
dois protocolos \emph{isolados um do outro}: a ordem lex (M1--M2) e a
covariância sob a dobra de índice ($E_{\partial}$). Nem um nem o outro nomeia
$C$. O $D_{\mathrm{can}}$ de \code{fis:thm:troca-realizacao} já prova
$u(Da,Db)=u(a,b)$ na árvore; \emph{isso não é} $E_{\partial}$.

\medskip\noindent\textbf{M1--M2 --- ordem lex, isolada de $\partial$.}
Cadeia $A_0\prec B_0\prec C_0\prec D_0$. Realização: ordem lex em Word.
\[
M_1:\ x\prec_{\mathrm{lex}}x'\Rightarrow f(x,y)\preceq_{\mathrm{lex}}f(x',y)
\qquad
M_2:\ y\prec_{\mathrm{lex}}y'\Rightarrow f(x,y)\preceq_{\mathrm{lex}}f(x,y')
\]
\begin{center}\footnotesize
\begin{tabular}{@{}lcc@{}}
\toprule
 & lexMax & rectCell \\
\midrule
$M_1$ & $0/24$ & $10/24$ \\
$M_2$ & $0/24$ & $6/24$ \\
\bottomrule
\end{tabular}
\end{center}
\emph{Classificação}: a ordem lex restringe a classe; \textbf{não} promove
$\operatorname{lexMax}$ a $C$.

\medskip\noindent\textbf{$E_{\partial}$ --- covariância da operação, isolada da lex.}
Alfabeto $(P_0,P_1,P_2,P_3)$, distinto da cadeia lex. A dobra no índice é
$D_{\mathrm{can}}(i)=(N-1)-i$. A pergunta \emph{não} é a isometria de $u$:
\[
E_{\partial}:\quad
\partial\bigl(f(\partial x,\partial y)\bigr)\;\stackrel{?}{\sim}\;f(x,y).
\]
lexMax: aplicável, \textbf{não} covariante ($4/16$ iguais, só a diagonal).
rectCell: imagem fora do carrier --- \textbf{não aplicável}.
E0 confirma $\partial^{2}=\mathrm{id}$ no alfabeto; isso \emph{não} determina
$R^{\dagger}$ nem $C$.

\medskip\noindent\textbf{Tuplo registado.}
\[
\boxed{
\begin{array}{c|cc}
 & \operatorname{lexMax} & \operatorname{rectCell} \\
\hline
M_1 & \checkmark & \times \\
M_2 & \checkmark & \times \\
E_{\partial} & \times\ \text{(aplicável, não covariante)} & \times\ \text{(não aplicável)} \\
\end{array}}
\]
\[
\boxed{\text{a ordem lex restringe; a dobra não selecciona; nenhuma determina }C.}
\]
Independência: $M_1\wedge M_2$ em lexMax $\not\Rightarrow E_{\partial}$.
\emph{Não se identifica} $E_{\partial}$ com a isometria $u(Da,Db)=u(a,b)$, nem
$D_{\mathrm{can}}$ com $\tau$, nem K6 com qualquer dos dois.

% ═══════════════════════════════════════════════════════════════════════════
\section{Síntese congelada}\label{comb:sintese}

\noindent\textbf{O que este ciclo fechou.}
\begin{enumerate}\itemsep3pt
\item Eliminação de falsas equivalências (C0--C8);
\item Capacidade determina preservação, não combinação (S0--S4);
\item Redução $\neq$ overwrite acidental (R0--R4);
\item Existência de combinação bilateral (D0--D3);
\item Não-unicidade estrutural (K0--K7);
\item Classe $\mathcal A_{\min}$ e invariantes discriminantes (L0--L7);
\item $\zeta/\mu/S$ compatível com ambas as geometrias --- não selecciona $f$;
\item Ordem lex restringe, $E_{\partial}$ não selecciona --- nenhuma determina $C$
(§\ref{comb:me}).
\end{enumerate}

\medskip\noindent\textbf{O que permanece aberto.}
\begin{enumerate}\itemsep3pt
\item Qual estrutura \emph{dentro da arquitectura} tem autoridade para seleccionar uma lei;
\item Operação $C\colon X\times X\to X$ \textbf{nomeada} e única;
\item Comparação com Hebb/Hopfield (§\ref{comb:hebb}) --- P6 \textbf{TRAVADA}.
\end{enumerate}

\medskip\noindent\textbf{Pergunta que substitui ``qual é $C$?''}
\caixa{%
  o que, dentro da arquitectura, teria autoridade suficiente para seleccionar uma lei?%
}

% ═══════════════════════════════════════════════════════════════════════════
\section{O que C0--C8 demonstrou (detalhe)}\label{comb:c0c8}

\noindent\textbf{Síntese C0--C8.}
\begin{enumerate}\itemsep3pt
\item $\text{overwrite}\neq\text{cópia}\neq\text{acumulação}\neq\text{combinação}$.
\item Fronteira $G_{\mathrm{state}}^{\mathrm{surv}}$: acumulação ($=2$) vs slot único ($=1$).
\item Comutatividade binária como estado; não-associatividade ternária no suporte cap~2.
\item Colisão mínima caracterizada.
\end{enumerate}

% ═══════════════════════════════════════════════════════════════════════════
\section{O que ainda não foi definido}\label{comb:aberto}

\noindent\textbf{Aberto --- após M1--M2 e $E_{\partial}$.}
\begin{enumerate}\itemsep3pt
\item Autoridade arquitectural suficiente para \textbf{seleccionar} uma lei entre
$\mathcal A_{\min}$ (ou extensões);
\item Nomear $C\colon X\times X\to X$ como lei única --- \textbf{não feito};
\item Comparação com regra Hebb / descida Hopfield (§\ref{comb:hebb}) --- P6 travada.
\end{enumerate}

\medskip\noindent\textbf{Explicitamente fechado.}
\begin{enumerate}\itemsep3pt
\item Existência de $f$ bilateral e bem definido;
\item Não-unicidade dentro de $\mathcal A_{\min}$;
\item Insuficiência de $\zeta+\mu$ para implicar $C$;
\item Insuficiência de impor L2--L4 apenas para ``fazer lexMax vencer'' --- isso seria
procurar a lei desejada; aqui registra-se \textbf{classificação};
\item Insuficiência de $M_1\wedge M_2$ e de $E_{\partial}$ para seleccionar $C$
(§\ref{comb:me}).
\end{enumerate}

% ═══════════════════════════════════════════════════════════════════════════
\section{Reserva para $C(x,y)$}\label{comb:reserva}

\begin{hipotese}[classe admissível --- medida, não escolha]\label{comb:hip-exist}
Existe uma classe
\[
\mathcal A_{\min}\subset\{f\colon X\times X\to X\}
\]
de operações observáveis bilateralmente dependentes, bem definidas, K6-invariantes,
e compatíveis com $G$ e $\zeta/\mu$. Dentro dela há pelo menos duas geometrias (lexMax, rectCell).
Impor comutatividade, idempotência e associatividade restringe a classe; \textbf{não}
identifica $C$ de forma única sem invariantes adicionais com autoridade arquitectural.
\end{hipotese}

\medskip\noindent\textbf{Estatuto.} Resultado \textbf{experimental} (D0--L7), não teorema de
$f=C$. lexMax é \textbf{sobrevivente} sob semilattice; rectCell é \textbf{contraexemplo}
à suficiência da dependência bilateral.

% ═══════════════════════════════════════════════════════════════════════════
\section{Relação futura com Hebb/Hopfield}\label{comb:hebb}

\noindent\textbf{Reserva explícita.} Este paper \textbf{não} é sobre Hopfield,
redes neurais clássicas nem regra Hebb. Regista o momento em que a arquitectura
multifocal chega à \textbf{necessidade} de uma operação que ainda \textbf{não possui}.

\medskip\noindent\textbf{Ponte futura} (quando $C$ existir como lei medida):
\begin{enumerate}\itemsep3pt
\item Comparar $C$ observada com $W_{ij}=\sum_p\xi_i^p\xi_j^p$
(\redes{Hopfield futuro});
\item Medir se $G_{\mathrm{real}}>1$ emerge \emph{depois} de $C$, não antes;
\item Manter P6 \textbf{TRAVADA} até a lei $C$ estar definida experimentalmente.
\end{enumerate}

\medskip\noindent\textbf{Escada completa (estado final).}
{\small
\[
\begin{array}{c}
\text{C0--C8}\ \checkmark
\ \to\ \text{S0--S4}\ \checkmark
\ \to\ \text{R0--R4}\ \checkmark
\ \to\ \text{D0--D3}\ \checkmark
\ \to\ \text{K0--K7}\ \checkmark
\ \to\ \text{L0--L7}\ \checkmark
\ \to\ \text{M1--M2}\ \checkmark
\ \to\ E_{\partial}\ \checkmark\\
\downarrow\\
\mathcal A_{\min}\ \text{medida}\ \to\ C(x,y)\ \textbf{reservada}\ \to\ \text{Hebb/P6 travada}
\end{array}
\]}
Este documento fecha na \textbf{classificação}; a lei única permanece reservada.

% ═══════════════════════════════════════════════════════════════════════════
\medskip\noindent\rule{\textwidth}{0.4pt}

\medido{\textbf{Bateria C0--L7 + M1--M2 + $E_{\partial}$}
(\code{node tests/redes\_combinacao.js}):\\
231 asserts, 0 falhas; \code{lib/arena\_combinacao.mjs};\\
integrada em \code{tools/auditoria\_rg6.bat} após §RG6.\\
\code{redes.tex} permanece interface (§RG6 fechado); a régua da troca alinha-se a
\code{fis:thm:troca-realizacao}.\\
\code{combinacao.tex} fecha com classificação; $C=\operatorname{lexMax}$ \textbf{não} declarado.}

\end{document}
